% Draft: §4, Mod--Head, for the restructured paper (rewrite method agreed
% 17 Oct 2026; see the evidence repo's DECISIONS.md). Drafted from the old
% draft's §§3.1, 3.4, 4, 6.1, 7, 8.1 and 10, CGEL ch. 5 §§7.10, 9.3--9.5,
% Payne et al. (2007), and the evidence repo's corpus records
% (corpus/premodification/, corpus/genitive/, corpus/exclusion_tests/).
% Revised after the muse review (reviews/muse-sec4-20261017-180414/TRIAGE.md).
% CGEL pages checked against the PDF on 17 Oct: 377, 392, 410, 415--419,
% 558; Payne et al. 569, 576, 579--581, 588. Old-draft pages not rechecked:
% 393--394, 549--550 (footnote), for the verified-detail pass.
% Constructed examples, judgments confirmed by Brett on 17 Oct: her every
% move and the various options good, *the in the room people bad.
% Added 18 Oct (the accounting of pronouns and names, Brett: "show the cards
% whichever way they lie"): CGEL pp. 327, 423, 429 (with n. 46), 430,
% 520--521 and Payne et al. p. 584 checked against the PDFs; Matthews (2014)
% p. 102 checked against the text conversion (p. 10 cut after the muse
% review, reviews/muse-sec4-6-20261018-065227/TRIAGE.md).
% 18 Oct (Brett: "would a tree help"): Figure fig:modhead added after the
% paragraph on Payne et al.'s tree; (b) follows their drawing of Head-Mod in
% (12b), p. 580, with a line from the nominal for each function. Halm (2021) pp. 19,
% 29--30 taken from the taxonomic paper's citation (determinatives-as-nouns.tex
% l. 838), which argues for a secondary use over conversion; not rechecked.
% Constructed examples, judgments confirmed by Brett on 18 Oct (M2): *Lucky
% few! bad, Lucky survivors were airlifted out good, Young Kennedys still go
% into politics good, *Remaining few were the hangers-on bad, Volunteers were
% precious few good. Poor everybody / Lucky everyone: Brett's COCA search,
% 18 Oct (M3), record corpus/exclusion_tests/2026-10-18-*.md in the evidence repo.
%
% Labels it relies on: sec:analyses (§2), sec:det-head (§3), sec:once (§5),
% sec:fused-relatives (§7). Not yet \input into the main file.

\section{Modifiers as heads: Mod--Head}\label{sec:mod-head}

In Mod--Head fusion the head of a nominal is also its modifier. \textit{CGEL} uses it for words that function elsewhere as modifiers in a nominal: determinatives and adjectives, and also nouns such as \mention{nylon} (p.~417), taken up in Section~\ref{sec:analyses}:\footnote{Examples (\ref{ex:mh}a, b, d--f) are \textit{CGEL}'s \citep[416--419]{huddleston2002}; (\ref{ex:mh}c) is from NOW (\textit{National Observer}, 3 October 2022).}

\ea\label{ex:mh}
\ea \mention{This copy is defective but \bk{the other two} are fine.}
\ex \mention{I didn't like it myself, but I respect the views of \bk{the many who did}.}
\ex \mention{Are you one of \bk{the lucky few who found powdered infant formula}?}
\ex \mention{\bk{The rich} cannot enter the kingdom of Heaven.}
\ex \mention{How will the new system affect \bk{the very poor}?}
\ex \mention{There were two pieces left and Kim, as always, chose \bk{the larger}.}
\z\z

For \textit{CGEL}, \mention{two}, \mention{many}, \mention{few}, \mention{rich}, \mention{poor}, and \mention{larger} each modify the nominal and head it at once \citep[415--419]{huddleston2002}. \textcite[579--580]{Payne2007} draw the adjectival case for \mention{the very hard of hearing}: the AdjP is \enquote{head of Nom} and \enquote{modifier of that head}, and \mention{the} is the NP's determiner as usual.\footnote{Payne et al.'s tree labels the fused function Head-Mod, since the counterpart modifier follows the head (\mention{people very hard of hearing}); their tree for \mention{this once} has Mod-Head (p.~588). \textit{CGEL} allows fusion with a post-head modifier as well (p.~418). I write Mod--Head for both.} Unlike Det--Head, both functions belong to the same nominal (p.~588; Figure~\ref{fig:modhead}). Their main Mod--Head case, \mention{this once}, is taken up in Section~\ref{sec:once}. On ordinary Head the phrase is the head of the nominal and nothing more.

\begin{figure}[H]
\centering
\begin{minipage}[t]{.36\linewidth}\centering
(a) \mention{the rich people}\par\smallskip
\begin{forest} nominal tree
[NP
 [{\synnode{Det}{DP}}
  [{\synnode{Head}{D}}, head edge [\mention{the}]]]
 [{\synnode{Head}{Nom}}, head edge
  [{\synnode{Mod}{AdjP}}
   [{\synnode{Head}{Adj}}, head edge [\mention{rich}]]]
  [{\synnode{Head}{N}}, head edge [\mention{people}]]]]
\end{forest}
\end{minipage}\hfill
\begin{minipage}[t]{.3\linewidth}\centering
(b) \mention{the rich}, fusion\par\smallskip
\begin{forest} nominal tree
[NP
 [{\synnode{Det}{DP}}
  [{\synnode{Head}{D}}, head edge [\mention{the}]]]
 [{\synnode{Head}{Nom}}, head edge
  [{\synnode{Mod--Head}{AdjP}}, no edge,
   tikz={\draw[-] (!u.south) -- ([xshift=-7pt].north); \draw[-,line width=1pt] (!u.south) -- ([xshift=7pt].north);}
   [{\synnode{Head}{Adj}}, head edge [\mention{rich}]]]]]
\end{forest}
\end{minipage}\hfill
\begin{minipage}[t]{.3\linewidth}\centering
(c) \mention{the rich}, ordinary Head\par\smallskip
\begin{forest} nominal tree
[NP
 [{\synnode{Det}{DP}}
  [{\synnode{Head}{D}}, head edge [\mention{the}]]]
 [{\synnode{Head}{Nom}}, head edge
  [{\synnode{Head}{AdjP}}, head edge
   [{\synnode{Head}{Adj}}, head edge [\mention{rich}]]]]]
\end{forest}
\end{minipage}
\caption{\mention{The rich people} (a), and \mention{the rich} under fusion (b) and ordinary Head (c). In (b) the AdjP has one parent, the nominal, with a line from it for each of its two functions; in Figure~\ref{fig:analyses}b the fused DP has two parents. In (c) the AdjP is the head of the nominal and nothing else.}\label{fig:modhead}
\end{figure}

So the two analyses agree that the adjective or determinative heads the nominal. Neither has a noun as the NP's ultimate head, the change of category Payne et al.\ accept at the fused node (p.~569). Both keep the word's category. \mention{Rich} and \mention{poor} take adverbs (\mention{the extremely rich}, \mention{the very poor}), which \textit{CGEL} calls \enquote{clear evidence that they retain their category status as adjectives}, and \mention{the rich} is plural with no plural ending (p.~418). Both give the head the dependents it takes elsewhere: a partitive in \mention{the biggest of them} (p.~416), a post-head phrase in \mention{the pure in heart} (p.~418). And neither restores a noun: \mention{the people very hard of hearing} can't be read generically, as \mention{the very hard of hearing} is \citep[580]{Payne2007}.

What's at issue is only the second function, the modifier. Form could show it in two ways: a form restricted to the fused function, as \mention{none} is for Det--Head, or a fact of order that only a modifier relation predicts. The formal differences Payne et al.\ cite are in Det--Head and in fused relatives (\mention{none}, p.~581; the \mention{-ever} forms, p.~576), and they cite none for Mod--Head. Their fact of order is adjacency: nothing comes between the fused functions, so \ungram{\mention{the in the room very hard of hearing}} is out (p.~580 n.~17). That argument is aimed at a silent head. With the AdjP as head, the PP would precede the head of a nominal, which PPs don't do before a noun either (\ungram{\mention{the in the room people}}).

For determinatives, \textit{CGEL} links the fused use to the dependent one. Determinatives \enquote{which can function as internal modifier as well as determiner can in general occur as fused modifier-head as well as determiner-head} \citep[415]{huddleston2002}: \mention{the few} beside \mention{the few mistakes}, where \mention{few} is modifier (p.~392). \textcite[588]{Payne2007} state the link without the hedge: \mention{some}, \mention{every}, \mention{any}, and \mention{no} \enquote{can function solely as determiners and therefore never permit a preceding determiner}.

Yet \mention{some} is attested after \mention{the}: \mention{when times are bad, \bk{the some} will become the many}; \mention{one of \bk{the some who hope for a return to the good old days}}.\footnote{COCA: \textit{Foreign Affairs}, 1999; iWeb: sspnet.org. COCA has five clear lines across searches for \texttt{the some}, \texttt{the some who}, and \texttt{the some of us who}; the other hits on the first display page of \texttt{the some} are typos for \mention{same}, the adverb \mention{some} \enquote*{approximately}, titles, and garbled text. iWeb has 10 hits for \texttt{the some who}, most with no earlier \mention{some} in view, so they aren't quotations of the word.} With \mention{the} as determiner, a fused \mention{some} can only be a modifier-head, and that's the Mod--Head Payne et al.'s \enquote{never} excludes. \textit{CGEL}'s licence is hedged, and it records exceptions of both kinds itself. \mention{Both}, which \enquote{cannot occur in internal modifier function}, heads a nominal after \mention{the} in \mention{the both of us} (p.~377). \mention{Every} and \mention{various} function as modifiers after a determiner (\mention{her every move}, \mention{the various options}) but don't head a nominal there: \enquote{\mention{Every} is again excluded, and so too this time is \mention{various}} (p.~416).

So the link between the two uses is a default, not a condition. A determinative that functions as modifier after a determiner can usually head a nominal after one too, but the exceptions run both ways, and both analyses list the determinatives that can. Ordinary Head states the list as a permission, as it states the list of determinatives that stand alone (Section~\ref{sec:det-head}). Both also need \textit{CGEL}'s observation that, cardinals apart, these determinatives are \enquote{hardly admissible} after a determiner unless the NP has a post-head modifier, as in \mention{these few here} and \mention{the many who did} (p.~416).

\mention{The some who} also tells against a silent noun (Section~\ref{sec:analyses}). In \mention{the some \textup{[N]} who say it}, a silent noun would make \mention{some} a modifier after \mention{the}, which Payne et al.\ exclude (p.~588), and the overt sequence doesn't turn up: of 53 iWeb hits for \mention{the some people}, none of the 50 displayed has that reading. The silent noun would license silently what's excluded overtly.

The determiners these heads take follow the word too. With an adjective before independent \mention{few}, the determiner is \mention{the} in 192 of 321 COCA lines, and a demonstrative (\mention{those chosen few}), a genitive (\mention{Hollywood's chosen few}), \mention{some}, \mention{we}, or \mention{us} in 35; nearly all the rest have \mention{a}, taken up below.\footnote{COCA, searched by hand for \mention{lucky}, \mention{chosen}, \mention{fortunate}, \mention{remaining}, and \mention{priveledged} (sic) before \mention{few}: 607 tokens, 368 lines inspected, 321 with no following noun. Of the 321, 192 have \mention{the}, 90 \mention{a} alone, 25 a demonstrative, 6 a genitive, 2 \mention{some}, and 2 \mention{we} or \mention{us}; one is a name, and three have nothing before the adjective (a headline, a fragment, and a web post).} That's the range of determinatives a plural noun takes. The adjectives of the \mention{rich} type take \mention{the} or a genitive (\mention{the world's poor}, \mention{our elderly}), though \textit{CGEL} says they're \enquote{determined by the definite article} and adds that \enquote{we couldn't even substitute a demonstrative: \ungram{\mention{these very poor}}} \citep[417]{huddleston2002}.\footnote{COCA, 18 October 2026, searched by hand for \texttt{[apps*] poor|rich|elderly}: of the 500 lines displayed, 19 have the adjective heading the NP after a genitive pronoun (\mention{care for \bk{their elderly}}, \textit{New York Times}, 1992), and 9 more quote \enquote{your tired, your poor}. Of 57 lines with a full stop after the adjective, 41 are the construction (\mention{employment for \bk{their poor}}, \textit{New York Times}, 1996), two of them with \mention{rich} (\mention{he doesn't mean \bk{our rich}}, reason.com, 2012). Demonstratives weren't searched.} Nor does the category fix the range. Among adjectives, \mention{blue} heads a nominal with no determiner (\mention{I like \bk{blue}}), \mention{second} heads one after \mention{a} (\mention{I didn't want \bk{a second}}), and \mention{bigger} can't (\ungram{\mention{a bigger}}) (pp.~416--417). Both analyses state the range word by word, or type by type.

What ordinary Head adds is one structure for every use of \mention{few}: \mention{few would dare}, \mention{few of them}, \mention{the few}, and \mention{the lucky few}. In each, \mention{few} heads the nominal, and a determiner, when there's one, is the NP's own. Fusion gives the first two Det--Head and the last two Mod--Head, so the same word heading the same kind of NP gets a different pair of functions depending on whether a determiner precedes it.

Which adjectives can head a nominal is a list on both analyses. \textit{CGEL} gives types: superlatives and definite comparatives (\mention{the tallest}, \mention{the elder}); ordinals (\mention{the second}); colour, provenance, and composition (\mention{I like \bk{blue}}); age and size (\mention{I prefer \bk{old}}); and the special interpretations of \mention{the rich} and \mention{the immoral} (pp.~416--417).\footnote{\textit{CGEL} treats one more case elsewhere. \mention{Above} is \enquote{primarily a preposition}, but belongs by conversion to the adjective category in the sense of location in the discourse, and the adjective is a fused modifier-head in \mention{Her conclusion is the same as \bk{the above}} (p.~1555). \textit{CGEL}'s own prose uses it freely (\mention{All the above have transitive uses}, \mention{Some of the above are verbs as well as nouns}, p.~1690). Like the types listed, it heads the nominal on both analyses. \textit{CGEL} adds that \mention{below} has no such use and \enquote{belongs only to the preposition category} (pp.~1555--1556), but COCA has both uses: before a noun in 185 of 197 lines of \mention{the below} with a following noun (\mention{are given in \bk{the below table}}, \textit{Education}, 2011), and heading a nominal in 18 of 55 lines of \mention{the below} before a full stop or after \mention{of} (\mention{is not either of \bk{the below}}, \textit{Popular Mechanics}, 2014).} It excludes \ungram{\mention{the advancing}} (p.~415), \ungram{\mention{a bigger}}, and \ungram{\mention{the main}} (p.~416). \mention{Main} is attributive-only (p.~558), so modifier is its only function, and its exclusion is \textit{CGEL}'s evidence that \enquote{the possibility of fusion with the head must in the last resort be stated in syntactic terms} (p.~416). Ordinary Head states the same list, as permissions to head a nominal.
Generality is also \textit{CGEL}'s argument against traditional grammar, which calls \mention{several} a pronoun in \mention{\bk{several} look quite promising} and treats \mention{the larger} as elliptical. Against that split, \textit{CGEL} holds that \enquote{a unitary treatment is to be preferred on grounds of simplicity and generality} \citep[419]{huddleston2002}, and fusion provides one. Ordinary Head keeps the gain with a single function: each word heads its nominal. Fusion gives the two words different pairs, Det--Head and Mod--Head.

Case shows no modifier function either. \textit{CGEL}'s argument from the genitive concerns the determiner function (p.~420), and it states no case requirement for modifiers, so the check here is this paper's own: whether the genitive divides the fused modifier-heads as it divides the determiner-heads (Section~\ref{sec:det-head}). It does. \mention{The latter's} and \mention{the former's} are frequent, \mention{the poor's} occurs in edited prose (\mention{\bk{the poor's} material well-being has improved}, \textit{Newsweek}, 2010), \mention{the rich's} only on the web, and \mention{few} after a determiner almost never.\footnote{One COCA search for \mention{the latter}, \mention{the former}, \mention{the other}, \mention{the poor}, and \mention{the rich} followed by \mention{'s} displayed 500 lines, probably a cap (498 pasted). In those lines the genitives number 149 for \mention{the latter}, 31 for \mention{the former}, 8 for \mention{the poor} (five in magazines and academic writing), and 3 for \mention{the rich} (all web or blog). \textit{CGEL} has one in its own prose: in its account of the traditional definition, a preposition \enquote{expresses the latter's relation to another word} (p.~598). \mention{The other} has 253, set aside because its plural \mention{others} leaves a noun analysis open. \texttt{the few/many 's} gave no genitive, and \texttt{few 's} one after a determiner, \mention{\bk{the first few's} salty freshness} (\textit{Ploughshares}, 1996, fiction).} In these lines too, the division follows the word.

Mod--Head is also where the advantage Section~\ref{sec:det-head} conceded to fusion has to be weighed. A bare determinative takes no pre-head adjective outside a few exclamations, taken up below (\ungram{\mention{Lucky few escaped}}),\footnote{\mention{Precious few would dare} isn't an exception. For \mention{precious} in \mention{They had precious little money}, \textcite[102]{matthews2014} offers the predicative use as the crucial evidence that it's an adverb (\mention{The money they had left was precious little}), and the same test works for \mention{few} (\mention{Volunteers were precious few}). So \mention{precious} is a modifier in the DP, as \mention{very} is in \mention{very few}, not in the nominal.} and fusion derives the gap from adjacency: with \mention{few} as Det--Head, \mention{lucky} would have to precede the determiner. After \mention{the}, \mention{few} is Mod--Head and \mention{lucky} modifies the nominal it heads. Ordinary Head states a condition: a determinative heading a nominal takes a pre-head adjective only when the NP has a determiner, or \mention{a} (taken up next).

\mention{A lucky few} is the hard case for both. Of the 321 COCA lines with an adjective before independent \mention{few}, 90 have \mention{a} as the only word before the adjective, and where a verb agrees with the phrase, it's plural: \mention{\bk{a lucky few} were sent far and wide to schools in other counties and states} (\textit{Washington Post}, 2015). The adjective can be a phrase (\mention{A very lucky few}, CNN, 2005). For fusion, \mention{a few} is \textit{CGEL}'s complex determinative, whose internal expansions are \mention{very}, \mention{good}, and \mention{fair} (p.~392). Fusion can open a slot for an AdjP inside that determinative, a phrase inside a word. It can make \mention{a} the NP's determiner, against \mention{a}'s requirement of a singular count head, since \mention{few} is plural. Or it can make \mention{few} a Mod--Head with no determiner, like \mention{blue} in \mention{I like \bk{blue}}, but in \textit{CGEL}'s analysis \mention{few} functions as modifier only after a determiner, and excluding \ungram{\mention{Lucky few escaped}} would then take the condition ordinary Head states, not adjacency.

On ordinary Head, \mention{a} in \mention{a few} isn't the determiner \mention{a}, which needs a singular count head. It's a modifier of the whole NP, the kind called peripheral in Section~\ref{sec:det-head}, though \textit{CGEL}'s peripheral modifiers are adverbs, PPs, and reflexive pronouns (p.~331): it comes after \mention{quite} (\mention{quite a few mistakes}) and before \mention{very} (\mention{a very few mistakes}; p.~392).\footnote{The expansions fit a phrase too: \textit{CGEL}'s \mention{a good few} and \mention{a fair few} (p.~392), and \mention{a select few}, with 21,137 tokens in NOW. For \mention{a great many}, \textit{CGEL} allows \mention{great} or \mention{good} and says that \enquote{one or other of these adjectives is required} (p.~394), but COCA has ten others in that slot, one token each, as in \mention{when a shrinking few do very well and \bk{a growing many} barely make it} (NPR, 2013). The slot takes an adjective before a noun with \mention{few} too: \mention{it took a few seconds, only a few seconds, \bk{a whole few seconds}, to sink in} (\textit{Jezebel}, 2016). \mention{A few} is positive where \mention{few} is negative (pp.~393--394), which I take to be a convention of the combination, as the minimizer reading of \mention{a bit} is in \mention{They didn't enjoy it \bk{a bit}} (pp.~549--550). \textit{CGEL}'s \mention{many a} (p.~394) comes apart the same way, as \mention{many} in predeterminer function before an NP with \mention{a}, like \mention{such} in \mention{such a mess}.} The condition on pre-head adjectives then has to count it (Figure~\ref{fig:afew}). So \mention{a lucky few} costs each analysis a statement: an AdjP slot inside a determinative, or an \mention{a} that admits the adjective.

\begin{figure}[H]
\centering
\begin{minipage}[t]{.24\linewidth}\centering
(a) \mention{a few}, fusion\par\smallskip
\begin{forest} nominal tree
[NP
 [{\synnode{Head}{Nom}}, head edge, before drawing tree={x+=1.5em}
  [{\synnode{Det--Head}{DP}}, no edge
   [{\synnode{Head}{D}}, head edge [\mention{a few}]]]
  {\draw[-] (!uu.south) -- (); \draw[-,line width=1pt] (!u.south) -- (); }]]
\end{forest}
\end{minipage}\hfill
\begin{minipage}[t]{.32\linewidth}\centering
(b) \mention{a few}, ordinary Head\par\smallskip
\begin{forest} nominal tree
[NP
 [{\synnode{Periph.\ Mod}{DP}}
  [{\synnode{Head}{D}}, head edge [\mention{a}]]]
 [{\synnode{Head}{NP}}, head edge
  [{\synnode{Head}{Nom}}, head edge
   [{\synnode{Head}{DP}}, head edge
    [{\synnode{Head}{D}}, head edge [\mention{few}]]]]]]
\end{forest}
\end{minipage}\hfill
\begin{minipage}[t]{.4\linewidth}\centering
(c) \mention{a lucky few}, ordinary Head\par\smallskip
\begin{forest} nominal tree
[NP
 [{\synnode{Periph.\ Mod}{DP}}
  [{\synnode{Head}{D}}, head edge [\mention{a}]]]
 [{\synnode{Head}{NP}}, head edge
  [{\synnode{Head}{Nom}}, head edge
   [{\synnode{Mod}{AdjP}}
    [{\synnode{Head}{Adj}}, head edge [\mention{lucky}]]]
   [{\synnode{Head}{DP}}, head edge
    [{\synnode{Head}{D}}, head edge [\mention{few}]]]]]]
\end{forest}
\end{minipage}
\caption{\mention{A few} under fusion (a), as \textit{CGEL}'s complex determinative, and under ordinary Head (b), and \mention{a lucky few} under ordinary Head (c). On ordinary Head \mention{a} is a peripheral modifier of the whole NP, attached as \textit{CGEL} attaches external modifiers (p.~331), and the NP it modifies has no determiner. In (c) the adjective is inside that NP's nominal, which is why ordinary Head's condition on pre-head adjectives has to mention \mention{a}.}\label{fig:afew}
\end{figure}

For \textcite[584]{Payne2007}, a condition of this kind is \enquote{precisely the kind of stipulation which is not required when the determiner and head functions are fused}. They say it of an analysis of the compounds with a silent head, but it applies to ordinary Head's condition as well. Whether the grammar needs such a statement anyway turns on the other heads that take no determiner in their primary use, pronouns and proper names. \textit{CGEL}'s treatment of them separates three questions: whether such a head takes a pre-head dependent at all, whether it takes a non-restrictive adjective, and whether a restrictive one needs a determiner.

On the first, the grammar already makes a statement of this kind, for pronouns. \textit{CGEL} makes it \enquote{a defining property of pronouns} that \enquote{they do not permit determiners} (p.~429), and they \enquote{do not normally allow internal pre-head dependents} (\ungram{\mention{Extravagant he bought a new car}}, p.~430). They take some dependents after the head, as in \mention{you who worked on both the projects} (p.~429). The compound determinatives split the same way (\mention{nothing sensible}, \mention{everything in the collection}), and for them \textit{CGEL} derives the gap from fusion: \enquote{Because determiner and head are fused it is not possible for these forms to take internal pre-head dependents} (p.~423). A pronoun has no determiner function to derive it from.

So the grammar already states, without deriving it, that one kind of head with no determiner normally takes no pre-head dependent. Ordinary Head's condition isn't that statement, though. Since \textit{CGEL} defines pronouns as taking no determiner, it can state their restriction without mentioning one. Pronouns with a determiner, as in \mention{the real me}, come up below. A determinative takes an adjective after a determiner or after \mention{a}, so ordinary Head's condition has to mention both. Fusion's derivation saves that statement, but not a kind of statement the grammar would otherwise lack.

The second question favours ordinary Head, though not by much. Names take a few non-restrictive adjectives with no determiner (\mention{poor Henry}, \mention{Young Isaac Newton went off to Cambridge}; pp.~520--521), and the core personal pronouns take \mention{lucky}, \mention{poor}, and \mention{silly} (\mention{Lucky you!}, \mention{poor old me}; p.~430). Simple determinatives don't seem to (\ungram{\mention{Lucky few!}}, \ungram{\mention{Poor many!}}). But two compounds are attested in the pronouns' exclamation, each just after the same exclamation with a name or a pronoun: \mention{Poor Val. Poor everybody.} (\textit{New Yorker}, 2005), \mention{Lucky you. Lucky everyone.} (\textit{Everwood}, 2004).\footnote{COCA, 18 October 2026, searched by hand for \texttt{poor|lucky|silly everybody|everyone|somebody|someone}: 28 lines. Of the rest, 19 have \mention{lucky} before a clause (\mention{You're lucky someone called it in}), four are predicative (\mention{How poor everyone had been}), one has a determiner (\mention{I made him the lucky someone}), and two are a list and a depictive. No line has \mention{somebody} or \mention{someone} in the use.}

For \textit{CGEL}, \mention{everybody} and \mention{everyone} are fused determiner-heads, so an internal pre-head dependent is \enquote{not possible} for them (p.~423). Fusion can call them nouns by conversion here, like \mention{a nobody} (p.~423 n.~43), though they keep the compounds' universal meaning and take no determiner. It can list them as an exception, as \textit{CGEL} does for pronouns, or take the adjective as a modifier of the whole NP, though that separates them from \mention{Lucky you!}, which \textit{CGEL} counts among internal pre-head dependents (p.~430). On any of these, adjacency no longer derives the gap for them. Ordinary Head adds them to \textit{CGEL}'s exception, \enquote{a few adjectives such as \mention{lucky}, \mention{poor}, \mention{silly} with the core personal pronouns}. Either way the compounds need a listed exception, and only on fusion does it go against a principle.

The third question is where pronouns and names look most like determinatives, but the likeness comes from countability. A restrictive adjective after a determiner is routine with personal pronouns: COCA has 350 tokens of \mention{the real me}. In this use the pronoun heads a singular NP, which takes a third-person singular verb (\mention{The real you is still alive back on Earth}, \textit{Analog}, 2019).\footnote{COCA, 15 October 2026, searches run by hand. Of the 14 lines for \texttt{the real you [be]}, ten have \mention{the real you} as subject, all with \mention{is} or \mention{'s} and none with \mention{are}. On one display page of \texttt{the real you}, 193 of the 198 lines have the restrictive use.} \textit{CGEL} would presumably call the word a common noun here, as it calls \mention{he} and \mention{she} in \mention{Is it a he or a she?} (p.~429 n.~46), though the reflexive suggests a secondary use of the pronoun.\footnote{The reflexive keeps the pronoun's person (\mention{The 12-year-old me is kicking myself}), as a common noun's doesn't (\ungram{\mention{Daddy is enjoying myself}}), except that a stage of the speaker contrasted with another can take third-person forms \citep[19, 29--30]{halm2021virtual}.} On either view the head is singular and count, and an NP with a singular count head needs a determiner, with or without a modifier.

Names are the same. In \mention{The young Isaac Newton showed no signs of genius}, \mention{young} restricts a secondary use in which the name denotes a set of manifestations, and in such uses \enquote{determiners are thus selected in the ordinary way} (pp.~520--521). A singular name in that use needs a determiner for the reason a singular count noun generally does. A plural head needs none, and takes a restrictive adjective without one, a plural name included (\mention{Lucky survivors were airlifted out}, \mention{Young Kennedys still go into politics}). So nothing that holds for names or for \mention{the real me} predicts \ungram{\mention{Remaining few were the hangers-on}} beside \mention{The remaining few were the hangers-on} (\textit{Christian Century}, 1998). Fusion's adjacency does.

The comparison cuts both ways, then. On the first question, \textit{CGEL} already states that pronouns normally take no pre-head dependent, so ordinary Head's condition is a further statement of a kind the grammar already contains. On the second, two compounds take an adjective in the pronouns' exclamation, which adjacency rules out. On the third, the facts with a determiner follow from number and countability and give ordinary Head nothing. The pre-head adjective remains the nearest thing to a trace of the determiner function (Section~\ref{sec:det-head}), but its derivation saves less than Payne et al.\ suggest.

Nothing in the form of Mod--Head expressions, then, shows the modifier function. Which determinatives take an external determiner, which determiners they take, which adjectives head a nominal, and which heads take the genitive all follow the word, and both analyses list them. Fusion derives the absence of an adjective before a bare determinative; after \mention{a}, and for \mention{Poor everybody}, both analyses need a statement.
